\chapter[Sermon 9. The Second Epistle of Jesus Christ by John to Them of Smyrna Is Expounded. And Is an Exhortation to Patience, and Consolation in Afflictions.]{The Second Epistle of Jesus Christ by John to Them of Smyrna Is Expounded. And Is an Exhortation to Patience, and Consolation in Afflictions.}
\label{sermon:009}
\markright{Sermon 9. The Second Epistle of Jesus Christ by John to Them of ...}

And unto the angel of the congregation of Smyrna write. These things saith he that is first and last, which was dead and is alive. I know thy works, and thy labors and poverty, but thou art rich. And I know the blasphemy of them which call themselves Jews, and are not: but are the congregation of Satan. Fear none of those things which thou shalt suffer. Behold, the Devil shall cast some of you into prison to tempt you, and you shall have tribulation for ten days. Be faithful unto death, and I will give thee a crown of life. Let him that hath ears, hear what the Spirit saith to the congregations; he that overcometh shall not be hurt of the second death.

Jesus Christ, from the right hand of the Father, through the ministry of an angel by the Apostle and Evangelist John, exhorts the congregations of Smyrna, who are afflicted with every kind of evil for the word of God, to endurance and comfort while they now suffer under the cross, promising great things to those who overcome. Truly, there can be no exhortation or consolation of this manner or in this matter that is briefer. For it is wisely couched in the eternal wisdom of the Father, and is applicable to all times, and it agrees very well with all who mourn under the cross. For just as Christ at the right hand of the Father is the catholic or universal Bishop, so too is his doctrine general, which he himself also applies to all congregations in the end of this Epistle and in others. And he declares that he loves his church, and is present in it by his power and aid.

And truly it is to be marveled at that nothing is condemned in this church, since some fault is found in conduct with all others. Therefore, the church of Smyrna was exceedingly excellent, though not without imperfections. For the Lord, in his goodness, does not hold small faults against us—of which the Prophet speaks, who will say, “My heart is clean? And from my hidden sins cleanse me”—so long as there is a fervent desire and zeal for godliness within us, and we are void of great evils.

\index{Apocalypse (Book of)}\index{John, Saint}\index{Jerome, Saint}\index{Irenaeus}\index{Eusebius of Caesarea}First, it is shown to whom this heavenly letter is sent, to the pastor of the church of Smyrna, and to the whole flock. For the captain is said to have sought peace, or fled, or taken peace, when the whole army together with him has done this. And the histories bear witness that Polycarp was that same messenger or pastor of the church of Smyrna, ordained by the Apostles themselves, namely by John the Apostle, Bishop there, and that he lived in the mystery of this congregation for eighty-six years. For so many years he declared himself before the Lieutenant Herod, when he was brought to execution. For in the fourth persecution of the church, Aurelius Antoninus and Aurelius Commodus, being Emperors, he was taken and brought before the governor. And at length, for the open and sincere confessing of Christ, he was burned. He had this very much in his mouth: that nothing ought to be received as true unless it were known to be set forth by the Apostles. Irenaeus affirms that when he was a child he saw this old father, a man of great years and reverence, in the third book and third chapter against heresies, where he tells many things of him besides. As also does Eusebius in the fourth book of the church history, in the fourteenth and fifteenth chapters. And Jerome in the register of the famous writers of the church. Eusebius in his Chronicles notes that he suffered martyrdom in the year of our Lord 190. Whereby it appears that he was made Bishop of Smyrna in the year of our Lord 84, or thereabout. For we said even now that he had been in that ministry for eighty-six years. And therefore, had he been Bishop of Smyrna many years before the setting forth of the Apocalypse, which was written in the ninety-seventh year, would God all pastors would set before their eyes this good Polycarp to be followed, of whom there remains a notable Epistle to the Philippians.

After again is the author of the Epistle declared, which is set forth with two titles, taken out of the former visions of John and description of Christ. Thus saith the first and the last. \&c. Whereby is signified the eternal divinity of Christ, which wants beginning and ending. And of himself is everlasting. There is added that he was dead, and lives again, that is to wit, has risen from the dead. And this beginning accords right well to the matter. For they perceive that whosoever are afflicted for Christ and his Gospel, have a Lord and patron more mighty and more faithful, which in no wise can be overcome. Who can also in death keep his, like as he raised Christ from the dead, to the intent we might have an opportunity, that we shall live with Christ, even in death it seems.

And now he comes to the matter itself, and that which he repeats in all his Epistles, he says here also: “I know your works, both good and evil. Do not think that I neither know nor care about your affairs. You are laid out in my hands; I know, see, and care for you and all yours.” These things greatly encourage us when we know that we have God watching over us, who has a care for us. And they also comfort us greatly, those who suffer, knowing that he who loves us and in no matter neglects us, has us always as it were before his eyes.

And here particularly he declares what he knew: First in deed the afflictions, which verily they suffered in the present persecution of the Emperor Domitian. And affliction is as it were a general word, encompassing four kinds of suffering. For he recounts, concerning their substance the spoiling of their goods and their poverty; in their name, insults, reproaches, or blasphemies; by imprisonment and bonds, yea and death also. For through these afflictions godly men are tested for the sake of the wicked. And in these may be included all other kinds of tribulation. The which the Epistle of Jesus Christ recites in a godly order. There is nothing therefore of these matters which the Lord Christ does not know.

Poverty has the first place. We should not take this to mean here a spiritual poverty, signifying humility of mind—though it is certain that the church of Smyrna possessed that virtue—but rather a poverty and lack of all things due to the confiscation of their goods. For in times of persecution, by virtue of royal proclamations, the goods of faithful Christians are confiscated for the king’s use, or permitted to soldiers, nobles, or others to take at their pleasure. The faithful, thrust out of their homes, are either driven into exile or reduced to begging. Would God we lacked examples of this today. Let us learn from this to bear and suffer such circumstances patiently, being persuaded that God knows our necessity. And because it is difficult for an honest man to hunger and want with his family, for comfort and consolation he adds, “But you are rich.”

This appears to the world as a paradox, or something unbelievable. What will they say, is he rich who has nothing and is reduced to the state of beggars? There are undoubtedly goods and riches of the mind much better than material possessions.

\index{Arethas of Caesarea}For this may be had, without the true felicity, of rich men of this world, who live a most miserable life. Again you shall see a poor man, concerning worldly goods, but furnished with the richness of the mind, for this cause only to be happy and most blessed. He covets nothing, he is content with his vocation: Neither would he change his state with the most wealthy and rich kings. Contrariwise you shall see rich men but of an evil conscience, and therefore thoughtful and burdened with cares, and never merry. You shall see poor men, but with merry hearts to lead a joyful life. Why then should it seem marvelous, if he that is stripped of his worldly goods for Christ, and enriched with the gifts of the mind, is glad and rejoices in God, and takes a good part in all chances, and for the same cause is judged to be verily rich? Doubtless the wise men of this world know also, that the only wise man is truly rich. Which is eloquently discoursed of Cicero. Arethas saith, in spiritual matters having a treasure hid in the field of thy heart, which is Christ, by reason of whom thou art rich also: Since thou hast him thy protector, who also when he was rich, for us became poor. \&c.

Secondly, blasphemy is mentioned, by which we understand all manner of slanderous remarks and insults, whereby the reputation and standing of the faithful are harmed. To this sort belong those who call the adherents of this faith heretics and schismatics; they deem them wicked people, despisers of God and his saints, and enemies of God’s service, and therefore plagues of the common wealth. If such people are tolerated, the common wealth must necessarily be destroyed. And these things often trouble good men more grievously than the loss of their possessions. For people value a good reputation as much as great wealth. Therefore, the Lord in the Gospel of Saint Matthew, chapter 10, heals this disease with many words and exhorts people not to do anything unworthy of the name of Christians to avoid such infamy.

\index{Antichrist}In the meantime, he declares also what motivated the authors of this mischief, whom he also blames severely, so that the faithful might understand how greatly God disapproves of the enemies of all godliness, and so that they might also care less for their hatred and persecution. They claim to be Jews, when they are nothing of the sort. Thus also Paul dealt with the Jews in his second letter to the Corinthians. The Jews are called confessors, honorers, and faithful servants of God. But these blaspheme God’s name, they attack the true faith, and oppress those who profess and worship God. Therefore, they are not Jews. What then? The synagogue, congregation, or assembly of Satan. Thus, the very Son of God plucks off the vermin from these young men, to the comfort of all those who suffer persecution. It should not grieve those who present themselves with proud titles that they are condemned by such harlots, the children of the Devil. Christ gives them a true title and calls them not the holy and catholic Church of God, but the conspiracy and school of Satan, as in whom, not the spirit of God, but of Satan, inspires lies, jugglings, deceits, blasphemies, fires, and deaths. Therefore, let it not grieve you at this day, even if it is your fortune to be condemned for the Gospel, by those who call themselves most holy, most shining, most reverent, and most irreproachable prelates and patrons of the old church, religion, and catholic faith, which have on their side, councils, fathers, so many successions of bishops, the prescription of so long a time, and the consent of so many realms. They desire nothing less than to be called such, but rather are the champions of Antichrist, and the professed enemies and tramplers underfoot of all Christian piety, for whom everlasting destruction is prepared.

After this, he offers a most evident exhortation and consolation, placing the Son before them, and saying, “Fear nothing of all that you shall suffer.” The Son of God himself feared the cross and death, and it is a natural thing to fear evils and death. Therefore, we are not commanded to be inhuman, nor to say, like stones, that the same things grieve us not, even though they torment us exceedingly. Rather, the faithful are encouraged to stand strong in the faith, and to do nothing unworthy of it for fear of punishment. We are therefore commanded boldly and freely to despise or suppress fear, and to seek strength by the Spirit of God, and to exercise it in temptations.

\index{Augustine, Saint}There follow reasons whereby he may obtain that which he hath persuaded, may confirm, comfort and exhort them with patience and constancy. He prophesies therefore to the godly, what things they shall suffer. And touches also the third kind of affliction, imprisonment and bonds, which we understand to mean all punishments whereby our bodies are tormented. But to be warned before of the evil is a great benefit. We are more easily overcome by unprovoked perils. And therefore the Lord in the Gospel after Augustine, the tenth chapter, and after John in the fifteenth and sixteenth chapters, tells his disciples of many evils that should come unto them, and adds thereto: These things have I spoken to you, that when the time shall come, you might remember that I have told you before. So now also faithfully warns the faithful in this Epistle.

And he touches the author of these evils, saying: The Devil will cast some of you into prison. Therefore we perceive that these evils arise from the common enemy of mankind and of the salvation of the faithful. Whereof we may conjecture that he seeks to intercept our salvation, and that we ought therefore to stand more earnestly against him. The soldiers when they hear that their old enemy is at hand, do not become sluggish, but cheerful. But the Devil inspires evil men, corrupts princes and magistrates, which attempt persecution against the church. We read that Satan afflicted Job, that is to have provoked the Chaldeans and Sabeans to kill his servants and drive away his cattle. Here therefore they may see with what authority they are encouraged, who at this day persecute the church of Christ, for the profession of the truth. The godly have that which may comfort them: For they hear that the same filthy beast is set against them, which so often has been vanquished by Christ the Prince of the faithful, and of the faithful through Christ's aid, may without any difficulty be overcome. And verily the Lord permits to the Devil and devilish men power over his servants. If you may well ask why, hear: That you may be tempted. God does not permit his people to Satan that they should perish, but that they should be tempted and tried. Therefore to a good end are delivered to the fire, that we might be purged from our sins, that the virtue of our faith might shine, and God might be glorified, and we made the purer. Who therefore will hereafter be impatient, when we hear that we for a great good are put to evil? We read in the third book of Wisdom: As gold is tried in the fire, so are the faithful proved. This parable has Saint Peter expounded at large in the fourth chapter of the first Epistle. Where he that will may have it more abundantly.

Moreover, the time of tribulation is also appointed, as it were, for ten days. The number ten signifies a multitude. For Jacob says to his father-in-law, “Ten times you have changed my wages.” Genesis 31. And Numbers 14 says he was tempted ten times, that is, often and many times. Job also affirms himself in chapter 29 to have been wronged ten times. Therefore, the Lord says at this present, “You shall be diversely and greatly troubled with evils.” Nevertheless, because he does not specify months, years, or ages, but days, he prophesies that the evils will not be continuous, but that there will always be spaces between to breathe in. Indeed, the shortness of persecution is first shown in Essay 26. Secondly, Saint Peter in 1 Peter 1 comforts the faithful. It is the part of the faithful not to prescribe to God, but whether we are put to pain for a long time or a short time, to take it patiently. Let us think rather that in the long continuance of evils, there is also some end foreseen by the Lord, and that in the same time of breathing, we must repair the evils and return to battle.

Finally, the godly are encouraged by a most ample and large promise, in which is also included the fourth and most grievous kind of affliction, and even bitter death itself, through fire, hanging, sword, water, etc. But if you are not afraid of death, but overcome it, and offer yourself to God, then I will give you, says the Lord, a crown of life. Hereunto is annexed the state of the Epistle, and some of all. Therefore be faithful, cheerful, constant, even to death. For the Lord also says in the Gospel: Whoever perseveres to the end shall be saved. And we read that the Apostle has said, if we die with Christ, we shall live with him. And truly, the crown of life is nothing other than eternal life, that everlasting, celestial, and unspeakable joy. And the Lord alludes to conflicts, after which, when fortunately finished, the victors are crowned. Blessed is the man, says the Apostle James, who suffers temptation, because when he is tried, he shall receive a crown of life, which the Lord has promised to those whom he loves. Similar things the Apostle Paul has also written in the first to the Corinthians, chapter nine, and in the second to Timothy, chapter four. Therefore let it be hard hereafter for no man to lose this temporal life, since, as it is lost for Christ, we shall receive eternal life; otherwise, we will, or will not, we must die. Let us therefore be content rather to die blessedly than to live miserably, so that we may please God.

\index{Aquinas, Thomas}Finally, just as in the end of the first Epistle, he communicated and applied the same wholly to all times and churches, lest any should suppose that these things concern them nothing. So in the end of this Epistle also, he both declares the Spirit to be the author of all these things, and exhorts all men to hear and obey diligently, and affirms it to be written unto all congregations in the world for edifying. Moreover, the promise of life he communicates, saying: “He who overcomes shall not be hurt by the second death.” This is spoken to all men and women, if you overcome. Therefore, we must overcome the world, the Devil, the flesh, and all temptation. And we must overcome by him, through faith, by his Spirit dwelling in us, and that we should walk in the way wherein he has commanded us to walk. If you overcome, you shall not be hurt in the second death. Thomas of Aquin says that the first death is of sin, the second of pain. We understand plainly by the first death the natural separation of the soul from the body, which also comes to us because of sin, as appears in the third chapter of Genesis. This same [illegible] has come to good and evil. For we are all earth, and to earth we shall return. And immediately follows the second death and the second life: Those who believe in Christ overcome, feel nothing of the second death, but live, as the Lord himself assures us in the third and fifth chapters of John. He shall not come into judgment, but has passed from death to life. But the wicked or unbelievers are conveyed straight ways from corporal death to everlasting death. Not that their souls can die, that is, cease to be, or that their bodies will not rise again, but that being deprived of that celestial and divine life of Christ, they feel everlasting torment, which state is rightly called death. These things are unknown to worldly men, who know no other life but this temporal. But God’s truth teaches us that there is another life and death after this, namely, the celestial and death infernal, or full of perpetual sorrows. That same is full of consolation, that we hear how the faithful, after paying the debt of this temporal life once, shall no more feel any torments. What then do the monks and friars prate of purgatory, babes? Let us praise our Savior Christ, who has delivered us from death and given us the hope of life everlasting, to whom be glory, praise, and so forth.
